\section{The high-\texorpdfstring{$\beta$}{beta} determinant-cell ledger}
\label{sec:tpc31-high-beta}

We insert the all-content theorems into the selected high-$\beta$
packet closed in TPC-28.  Its exact exponents are
\begin{equation}
 \beta=\frac{267}{400},
 \qquad
 \jmath=1-\beta=\frac{133}{400}.
 \label{eq:tpc31-high-parameters}
\end{equation}
Thus
\begin{equation}
 Q=X^{267/400+o(1)},
 \qquad
 J=X^{133/400+o(1)}.
 \label{eq:tpc31-high-QJ}
\end{equation}
Every conclusion in this section is an upper bound for a
support-compatible cell.  It does not assert that the cell is
nonempty or has positive asymptotic mass.

\subsection{A translated orbit-scale window}

Let $I_X$ be any interval satisfying
\begin{equation}
 |I_X|=X^{133/400+o(1)}=JX^{o(1)}.
 \label{eq:tpc31-high-J-window}
\end{equation}
The interval is initially allowed any center.  For the fully coprime
target layer, \cref{cor:tpc31-coprime-window} gives the relaxed margin
\begin{equation}
 \beta-\jmath
 =\frac{267}{400}-\frac{133}{400}
 =\frac{134}{400}
 =\frac{67}{200}.
 \label{eq:tpc31-high-coprime-J-window-margin}
\end{equation}

For all target contents, use
\eqref{eq:tpc31-all-window-saving}.  The large-content side limits the
splice to
\begin{equation}
 \min\left\{\frac{133}{400},\frac{134}{400}\right\}
 =\boxed{\frac{133}{400}}.
 \label{eq:tpc31-high-all-J-window-margin}
\end{equation}

\begin{corollary}[All-content orbit-scale translated window]
\label{cor:tpc31-high-J-window}
For every fixed $\sigma<133/400$, the sum of the three physical raw
channels restricted by
\begin{equation}
 \Delta^\#_{\alpha,\gamma}(j)\in I_X
 \label{eq:tpc31-high-J-window-event}
\end{equation}
satisfies
\begin{equation}
 \boxed{
 \sum_{\nu=1}^3
 |\cS_\nu(\Delta^\#\in I_X)|
 \ll_{\sigma,h,\mathscr D}XQX^{-\sigma}.}
 \label{eq:tpc31-high-J-window-bound}
\end{equation}
\end{corollary}

\begin{proof}
In \cref{cor:tpc31-optimized-cell}, take $q=1$ and
$s=133/400$.  Equation
\eqref{eq:tpc31-high-all-J-window-margin} is the resulting endpoint.
The $o(1)$ in \eqref{eq:tpc31-high-J-window} is absorbed by the strict
choice of $\sigma$.
\end{proof}

If the same window is genuinely macro-centered, the all-content result
improves because large content cannot occur.

\begin{corollary}[Macro-centered orbit-scale window]
\label{cor:tpc31-high-macro-J-window}
Assume in addition that, for a fixed $\lambda>0$,
\begin{equation}
 I_X\subset\{\delta:|\delta|\ge\lambda Q\}.
 \label{eq:tpc31-high-macro-J-condition}
\end{equation}
Then every fixed $\sigma<67/200$ satisfies
\begin{equation}
 \boxed{
 \sum_{\nu=1}^3|\cS_\nu(\Delta^\#\in I_X)|
 \ll_{\sigma,h,\lambda,\mathscr D}XQX^{-\sigma}.}
 \label{eq:tpc31-high-macro-J-bound}
\end{equation}
\end{corollary}

\begin{proof}
Use \cref{prop:tpc31-macro-content-collapse} with $q=1$ and support
exponent $s=133/400$.  The available margin is
$\beta-s=134/400=67/200$.
\end{proof}

\subsection{A wider macroscopic window}

Take instead
\begin{equation}
 |I_X|=X^{1/2+o(1)}
 \label{eq:tpc31-high-half-window}
\end{equation}
and again allow the center to be $\asymp Q$.  Since
\begin{equation}
 \beta-\frac12
 =\frac{267}{400}-\frac{200}{400}
 =\boxed{\frac{67}{400}},
 \label{eq:tpc31-high-half-margin}
\end{equation}
the cell remains power-small despite containing many more than $J$
normalized determinant values.

\begin{corollary}[All-content half-scale translated window]
\label{cor:tpc31-high-half-window}
For every fixed $\sigma<67/400$,
\begin{equation}
 \boxed{
 \sum_{\nu=1}^3
 |\cS_\nu(\Delta^\#\in I_X)|
 \ll_{\sigma,h,\mathscr D}XQX^{-\sigma}}
 \label{eq:tpc31-high-half-window-bound}
\end{equation}
whenever \eqref{eq:tpc31-high-half-window} holds.
\end{corollary}

\begin{proof}
Use \eqref{eq:tpc31-all-window-saving} with window exponent $1/2$:
\[
 \min\left\{\frac{133}{400},
             \frac{267}{400}-\frac12\right\}
 =\frac{67}{400}.
\]
\end{proof}

\subsection{One orbit-scale residue class}

Let $q_X$ be any integer sequence with
\begin{equation}
 q_X=X^{133/400+o(1)}.
 \label{eq:tpc31-high-residue-modulus}
\end{equation}
No primality or coprimality condition on $q_X$ is needed for the
canonical determinant cell.  For every choice of
$a_X\pmod{q_X}$, \cref{cor:tpc31-optimized-residue} gives
\begin{equation}
 \min\left\{\frac{133}{400},
             \frac{267}{400},
             \frac{133}{400}\right\}
 =\boxed{\frac{133}{400}}.
 \label{eq:tpc31-high-residue-margin}
\end{equation}

\begin{corollary}[All-content orbit-scale residue cell]
\label{cor:tpc31-high-residue}
For every fixed $\sigma<133/400$,
\begin{equation}
 \boxed{
 \sum_{\nu=1}^3
 |\cS_\nu(\Delta^\#\equiv a_X\!\pmod{q_X})|
 \ll_{\sigma,h,\mathscr D}XQX^{-\sigma}.}
 \label{eq:tpc31-high-residue-bound}
\end{equation}
\end{corollary}

The intersection of an interval of length $\asymp J$ with one class
modulo $q_X\asymp J$ contains $O(1)$ integers.  The general all-content
splice has margin $133/400$.  On the fully coprime layer alone the
margin is $\beta=267/400$; and if the interval is macro-centered as in
\eqref{eq:tpc31-high-macro-J-condition},
\cref{prop:tpc31-macro-content-collapse} shows that the same
$267/400$ margin holds for all contents.  Without that location
hypothesis, the stronger margin must not be quoted for the all-content
cell.

\subsection{Separation from earlier deletions}

The translated intervals above may be chosen inside
\begin{equation}
 \kappa Q\le|\Delta^\#|\le2\kappa Q
 \label{eq:tpc31-high-macro-center}
\end{equation}
for fixed $\kappa>0$, whenever the physical row support meets that
annulus.  Then \cref{prop:tpc31-macro-content-collapse} gives
$c\ll_{\kappa,\mathscr D}1$ and
$|m_\alpha-m_\gamma|\asymp Q$, so this lies outside a near-row
deletion.  On the fully coprime sublayer, every selected reflected gcd
equals one by \cref{prop:tpc31-selected-dilate}; the
large-full-content theorem of TPC-30 and the large-reflected-content
part of TPC-29 do not estimate that subcell.  Some expanded divisor subfamilies may overlap a
different TPC-29 sparse region, but the present bound is uniform for
the complete unexpanded raw kernel and does not assume that sparse
condition.

\begin{table}[H]
\centering
\caption{Margins supplied by the high-$\beta$ determinant-cell bounds.}
\label{tab:tpc31-high-ledger}
\begin{tabular}{@{}
 >{\raggedright\arraybackslash}p{0.46\textwidth}
 >{\centering\arraybackslash}p{0.17\textwidth}
 >{\raggedright\arraybackslash}p{0.26\textwidth}@{}}
\toprule
Cell & Available relaxed margin & Source \\
\midrule
Fully coprime $J$-window
& $67/200$ & row support count $J/Q$ \\
General all-content $J$-window
& $133/400$ & large-content splice \\
Macro-centered all-content $J$-window
& $67/200$ & bounded-content collapse \\
All-content $X^{1/2}$-window
& $67/400$ & window cardinality \\
All-content class modulo $J$
& $133/400$ & modulus/splice \\
General all-content intersection: $|I|\asymp q\asymp J$
& $133/400$ & large-content splice \\
Macro-centered all-content intersection
& $267/400$ & bounded-content one-cell floor \\
\bottomrule
\end{tabular}
\end{table}

Each table entry is the relaxed margin supplied by the stated bound,
not an optimality claim.  Every strictly smaller fixed exponent is
proved; no endpoint equality is claimed.
